只有主码保留候选才访问残差码。这使系统先对 $n_c$ 个候选进行广泛而便宜的扫描，再执行较窄的 $O(cd)$ 评分阶段。如果对每个访问候选读取残差，坐标计算量将趋向 $O(n_cd)$。lane 协作加速已选候选，不负责决定精排预算。

\subsection{重建评分与查询相关常数}

令 $\widehat p(x)$ 为主码重建，$\widehat r(x)$ 为变换坐标中的残差解码。编码器存储 $\gamma(x)=2\langle\widehat p(x),\widehat r(x)\rangle$。实现中的精排排序分数为
\begin{align}
 \refineddist(q,x)
 &=\primarydist(q,x)+\|q'-\widehat r(x)\|_2^2+\gamma(x)
 \label{eq:5-refined-score}\\
 &=\|q'-\widehat p(x)-\widehat r(x)\|_2^2+\|q'\|_2^2.
 \label{eq:5-refined-constant}
\end{align}
第二个等式由展开平方范数得到。因此，该分数按平方距离对层次重建排序，但返回值还包含一个仅与查询有关的常数；它既不是直接计算的原始向量距离，也不是未平移的重建距离。该常数不影响固定查询的排序。即使按此分数精确选择，重建误差仍然存在\cite{jhq}。

残差坐标使用子空间内共享的标量码本项。标量路径遍历每个候选的各维，累加与 $q'$ 的平方差。查询坐标及在容量允许时的残差码本会搬入共享内存，最后再加上已存的修正项和主码分数。

\subsection{Lane 到 word 的映射与部分和归约}

对 $\lrefine=L\in\{1,8,16,32\}$，独立的精排 block 为每个候选分配 $L$ 个相邻 lane。一个 32-bit 残差 word 包含 $d_{\rm word}=32/B_R$ 个坐标编码，其中 $B_R$ 为残差位宽。被测八比特路径中 $d_{\rm word}=4$，四比特路径中 $d_{\rm word}=8$。lane $a$ 读取 word 索引 $a,a+L,a+2L,\ldots$，解码坐标并累加部分和。butterfly shuffle 归约合并 $L$ 个部分和，lane zero 写出精排分数。

\begin{figure}[tbp]
\centering
% Lane mapping for L=16; schematic, not a measured trace.
\begin{tikzpicture}[font=\small,
 box/.style={draw=ThesisNavy!65,rounded corners=2pt,align=center,
   text width=5.8cm,minimum height=1.5cm,inner sep=5pt},
 arr/.style={-{Latex[length=2mm]},draw=ThesisNavy!80,thick}]
\node[box,fill=ThesisNavy!5] (a) at (0,0)
 {Warp lanes 0--15: candidate $x_0$\\
  subgroup lane $a$: words $a,a+16,\ldots$};
\node[box,fill=ThesisNavy!5] (b) at (7,0)
 {Warp lanes 16--31: candidate $x_1$\\
  subgroup lane $a$: words $a,a+16,\ldots$};
\node[box] (ar) at (0,-2.2)
 {Decode coordinates; accumulate partials\\
  Shuffle offsets $8,4,2,1$};
\node[box] (br) at (7,-2.2)
 {Decode coordinates; accumulate partials\\
  Shuffle offsets $8,4,2,1$};
\node[box,text width=11.5cm,minimum height=1.15cm] (out) at (3.5,-4.25)
 {Lane zero of each subgroup writes\\
  $\refineddist=\primarydist+\text{residual squared-difference sum}+\gamma(x)$};
\draw[arr] (a) -- (ar);
\draw[arr] (b) -- (br);
\draw[arr] (ar.south) -- (out.north west);
\draw[arr] (br.south) -- (out.north east);
\end{tikzpicture}

\caption[十六 lane 残差精排]{十六 lane 残差精排。一个 warp 处理两个候选；候选内的 lane $a$ 读取残差 word $a,a+16,\ldots$。残差编码为八比特时，每个 32-bit word 解码四个坐标。子组归约合并部分和，再加上保留的主码分数和修正项。即使子组内 word 访问连续，候选物理位置仍可能分散。}
\label{fig:5-lane-refinement}
\end{figure}

启动条件要求残差行按 word 对齐，且 $D_s$ 能被 $d_{\rm word}$ 整除，从而一个 word 不跨越两个子空间的码本边界。所有 lane 都执行 shuffle 序列，包括分配到无效或超出范围槽位的 lane；它们在输出标为无穷前贡献零。这避免让 warp 参与情况取决于某个候选是否存在。$L=16$ 时，一个 warp 处理两个候选，而不是每个 lane 各处理一个候选。

实验将两种改变分开。L0 是原来的标量逐字节加载路径，L1 使用 word load，但仍每候选一个 lane，保持其累加顺序。L8/L16/L32 进一步将坐标工作分配给多个 lane，缩短每个 lane 的串行遍历，并在 $c$ 较小时提供有用工作；代价是增加 shuffle，并减少同时处理的独立候选数量。因此 $L$ 并非越大越好。保留候选位置不规则、码本处于共享还是全局内存，也仍影响成本。

\subsection{最终选择与数值边界}

当至少有 $k$ 个有效候选具有有限精排分数、且数据库 ID 唯一时，最终 kernel 按 $(\refineddist,\mathrm{database\ ID})$ 精确选择 top-$k$。每条查询一个 block，读入 $c$ 个分数后反复归约最小值，输出其 ID 和分数，再移除该分数进入下一轮。每轮扫描全部 $c$ 个槽位，没有固定的逐线程候选配额。该契约要求至少 $k$ 个有限有效条目。退化池中，选择器将已选分数设为无穷，但保留有效位置：有限条目用完后，最小的有效 ID 可能以无穷分数重复输出。全无效池则返回无效 ID 和无穷分数。这些情形超出本文评测使用的有限候选池契约。

最终 ID 并列规则与主码保留使用的物理位置规则不同。精确的最终选择器不能找回路由或主码筛选丢掉的候选，也不能消除量化误差。多 lane 归约在选择前改变浮点结合顺序，因此即使实数公式不变，近似并列候选也可能发生排序变化。

\Cref{sec:7-3}区分这些影响：L1 在被测用例中逐位一致，多 lane 的距离相对误差最高约为 $1.6\times10^{-6}$。普通预算下 L16 端到端收益范围为 $1.015$--$1.789\times$；内核级收益更大，不能替代端到端数字。受控的 load-only 与协作变体支持设计结论，而数值等价性需要检查，不能预先假定。
